% =====================================================================
%  Lời mở đầu và bảng đối chiếu với sáu hạng mục của bài
% =====================================================================
\chapter*{Lời mở đầu}
\addcontentsline{toc}{chapter}{Lời mở đầu}
\markboth{Lời mở đầu}{Lời mở đầu}

Các bài trước làm việc với dữ liệu mà thứ tự không quan trọng: một bệnh nhân, một căn nhà, một bức ảnh. Assignment 06
chuyển sang dữ liệu chuỗi, nơi thứ tự chính là thông tin. Giá cổ phiếu hôm nay phụ thuộc hôm qua; một người sắp huỷ gói
nghe nhạc thường nghe thưa dần trước khi dừng hẳn. Mạng nơ-ron hồi quy được tạo ra cho đúng loại dữ liệu này: nó đọc
chuỗi từng bước và mang theo một trạng thái nhớ.

Bài dựng bốn kiến trúc hồi quy (Simple RNN, LSTM, GRU, BiLSTM) bằng cả PyTorch và Keras, huấn luyện trên hai tập dữ liệu
thật, tổng cộng 16 mô hình. Tập thứ nhất là giá cổ phiếu Amazon trong 27 năm, một bài hồi quy. Tập thứ hai là log nghe
nhạc theo ngày của \SL{kkbox/n_sample} người dùng KKBox, một bài phân loại người dùng rời bỏ dịch vụ. Chương 1 dựng lý
thuyết từ công thức, mỗi khái niệm là một hàm NumPy đối chiếu với PyTorch. Chương 2 phân tích và tiền xử lý hai tập.
Chương 3 và 4 huấn luyện 16 mô hình rồi so sánh hai framework. Chương 5 đóng gói các mô hình thành một Web App.

Hai tập cho hai kết quả trái ngược, và cả hai đều có ích. Trên KKBox, các mô hình học được một tín hiệu thật: người sắp
rời bỏ nghe nhiều hơn ở đầu tháng rồi giảm dần, và LSTM đạt ROC-AUC \SL{torch_kkbox/models/lstm/roc_auc} chỉ từ chuỗi
hành vi. Trên AMZN, không mô hình nào trong 16 mô hình thắng được mốc ngây thơ ``giá ngày mai bằng giá hôm nay''
(RMSE \SL{bench/naive/rmse} USD). Báo cáo trình bày kết quả này như nó là, kèm phân tích vì sao, thay vì dựa vào chỉ số
$R^2$ xấp xỉ 0,99 trên giá, một chỉ số trông rất đẹp nhưng mốc ngây thơ cũng đạt được.

Mọi con số trong báo cáo được sinh tự động từ tệp kết quả của notebook, mọi đoạn mã được trích thẳng từ mã nguồn, và mọi
hình được vẽ bằng TikZ từ dữ liệu, kể cả hai khung giao diện của Web App. Chạy lại notebook thì báo cáo đổi theo.
Phụ lục~\ref{ch:phuluc} ghi các lệnh để tái lập.

Em xin cảm ơn thầy \textbf{PGS.TS Trần Đình Quế} đã hướng dẫn học phần.

\vspace{8mm}
\begin{flushright}
\emph{Hà Nội, tháng 9 năm 2026}\par
\vspace{2mm}
\textbf{Nguyễn Duy Nghĩa}
\end{flushright}

\section*{Đối chiếu với các hạng mục của bài}

\begin{table}[H]
  \centering
  \small
  \begin{tabular}{L{4.4cm}L{1.3cm}L{4.2cm}L{4.2cm}}
    \toprule
    Hạng mục & Chương & Notebook & Mã nguồn \\
    \midrule
    1. Cơ sở lý thuyết RNN, kèm code & \ref{ch:lythuyet} & \texttt{00\_ly\_thuyet\_rnn} & \texttt{a06/theory.py} \\
    2. Phân tích và tiền xử lý hai tập dữ liệu chuỗi thực tế & \ref{ch:dulieu} & \texttt{01\_du\_lieu\_amzn}, \texttt{02\_du\_lieu\_kkbox} & \texttt{a06/data\_stock.py}, \texttt{a06/data\_churn.py} \\
    3. Hiện thực và huấn luyện bằng PyTorch & \ref{ch:pytorch} & \texttt{03\_pytorch\_amzn}, \texttt{04\_pytorch\_kkbox} & \texttt{a06/models\_torch.py}, \texttt{a06/train\_torch.py} \\
    4. Hiện thực Keras, đối chuẩn PyTorch và Keras & \ref{ch:keras} & \texttt{05\_keras}, \texttt{06\_doi\_chuan} & \texttt{a06/models\_keras.py}, \texttt{a06/bench.py} \\
    5. Đóng gói triển khai Web App & \ref{ch:webapp} & -- & \texttt{app/} \\
    6. Kết luận & \ref{ch:ketluan} & -- & -- \\
    \bottomrule
  \end{tabular}
\end{table}
